% =============================================================================
% 5. fejezet — Összefoglalás és jövőbeli munka
% -----------------------------------------------------------------------------
% Az eredményekre vonatkozó konkrét számértékek a végleges futtatásokból
% töltendők ki (a [kitöltendő] jelölésű helyeken).
% =============================================================================

% =============================================================================
\section{A dolgozat eredményeinek összegzése}
% =============================================================================

A dolgozatban egy végponttól végpontig megvalósított, AI-alapú
bőrelváltozás-elemző rendszert (DermAI) terveztünk és valósítottunk meg, amely a
korai bőrrákszűrést egy okostelefonon és webböngészőben egyaránt elérhető
eszközzel támogatja. A munka két, egymást kiegészítő pillérre épült.

A \textbf{kutatási rész} egy Vision Transformer modellt tanított be az ISIC 2024
adathalmazon a bőrelváltozások jó\-indulatú/rosszindulatú osztályozására,
megfelelő figyelmet fordítva az adathalmaz rendkívül erős
osztály-egyensúlytalanságára (súlyozott mintavételezés, fókuszveszteség,
kétfázisú tanítás). A modellt két viszonyítási alappal vetettük össze: egy azonos
feltételek mellett betanított ResNet-18 konvolúciós hálóval (az architektúra
hatásának izolálására) és egy zero-shot Gemini nagy nyelvi modellel. A
kiértékelés a hivatalos ISIC 2024 részleges AUC metrikán, két klinikailag
értelmezhető működési ponton, valamint statisztikai szignifikancia-vizsgálattal
(DeLong-teszt, konfidencia-intervallumok) és kalibrációs elemzéssel (Brier-pontszám)
történt. \textit{[A fő eredmények -- a ViT, a ResNet-18 és a Gemini AUC- és
pAUC-értékei, valamint a statisztikai összehasonlítás -- a végleges futtatások
alapján kerülnek be: kitöltendő.]}

A \textbf{mérnöki rész} a betanított modellt egy teljes, működő alkalmazássá
egészítette ki: natív Android kliens, webes felület, két backend (Flask, illetve
FastAPI), a modellt kiszolgáló MCP-réteg, valamint a Gemini nagy nyelvi modellre
épülő képvalidáció, magyarázatgenerálás és párbeszéd. A rendszer a felhasználói
adatokat hitelesített hozzáférés mögött, felhasználónként elkülönítve tárolja
(Firebase), a token-szintű válaszközvetítést szerver-küldött eseményekkel valósítja
meg, és több szintű hálózati hibatűréssel alkalmazkodik a változó környezethez.

% =============================================================================
\section{Korlátok}
% =============================================================================

Az eredmények értelmezésekor több korlátot figyelembe kell venni:

\begin{itemize}
  \item \textbf{Kevés pozitív eset.} Az ISIC 2024 adathalmazban mindössze 393
  rosszindulatú eset szerepel, a teszthalmazban 59. Ez a metrikák -- különösen
  az érzékenység és a pAUC -- becslését érzékennyé teszi a mintavételi zajra,
  amit a megadott konfidencia-intervallumok tükröznek.
  \item \textbf{Kizárólag képi bemenet.} A modell csak a képet használja; az ISIC
  2024 verseny tapasztalata szerint a klinikai metaadatok bevonása érdemben
  javíthatja a teljesítményt.
  \item \textbf{Adathalmaz-jellegű torzítások.} A 3D-TBP felvételek nem fedik le a
  hétköznapi okostelefonos fényképezés teljes változatosságát, és a populációs
  reprezentáció (például a bőrtónusok eloszlása) korlátozott lehet, ami az
  általánosíthatóságot befolyásolja.
  \item \textbf{Tájékoztató jelleg.} A rendszer nem klinikai diagnosztikai
  eszköz; kimenete nem helyettesíti a szakorvosi vizsgálatot.
\end{itemize}

% =============================================================================
\section{Jövőbeli fejlesztési irányok}
% =============================================================================

A dolgozat eredményei több irányban is folytathatók:

\begin{itemize}
  \item \textbf{Multimodális fúzió.} A képi modell és a klinikai metaadatok
  együttes felhasználása (például gradiens-alapú döntésifa-együttessel
  kombinálva), az ISIC 2024 győztes megoldásainak mintájára, várhatóan tovább
  javítaná a teljesítményt.
  \item \textbf{Kalibráció és bizonytalanságbecslés.} A kimeneti valószínűségek
  utólagos kalibrálása (például hőmérséklet-skálázással) és a predikciós
  bizonytalanság becslése növelné a klinikai megbízhatóságot.
  \item \textbf{Robusztusság és méltányosság.} Sokszínűbb -- különböző
  bőrtónusokat és felvételi körülményeket lefedő -- adathalmazokon történő
  tanítás és kiértékelés csökkentené a populációk közötti
  teljesítménykülönbségeket.
  \item \textbf{Magyarázhatóság.} A figyelmi térképek vagy más vizuális
  magyarázó technikák integrálása átláthatóbbá tenné a modell döntéseit a
  felhasználó és a szakorvos számára.
  \item \textbf{Klinikai validáció.} A rendszer valós környezetben, prospektív
  vizsgálatban történő értékelése elengedhetetlen a tényleges klinikai
  hasznosság igazolásához.
\end{itemize}

Összességében a dolgozat azt mutatta meg, hogy egy modern, figyelem-alapú képi
modell és egy nagy nyelvi modell összehangolt alkalmazásával egy hozzáférhető,
felhasználóbarát bőrrákszűrő rendszer építhető, amely -- a fenti korlátok
tudatában -- értékes első lépés lehet a korai felismerés támogatásában.
